%%% FIELD proof-finite-inversion-algorithm
Let \(r\) be the number of distinct values in the fixed policy grid.  A propensity has one of the \(r+1\) positions in the open intervals cut out by those values or is tied with one of the \(r\) values.  Thus it has at most \(2r+1\leq 2L+1\) weak positions.  Retaining only pairs consistent with the propensity order can only reduce their number, so the definitions \(\texttt{def:weak-order-cells}\), \(\texttt{def:common-rank-paths}\), and \(\texttt{def:component-index}\) give
\[
 |\Omega(p^{1:L})|\leq (2L+1)^2,
 \qquad
 |\Gamma_K|=|\Omega(p^{1:L})|\,|\Sigma_K|
 \leq (2L+1)^2\binom{2K-2}{K-1}.
 \tag{1}
\]
Here the path count is the count proved in \(\texttt{thm:path-representation}\).  Consequently all order--path labels can be enumerated in finitely many steps.

Write
\[
 \mathcal Q_n=R_n\cap\overline{\mathcal Q}.
\]
Suppose first that \(\mathcal Q_n\neq\varnothing\).  For a fixed \(\gamma=(\omega,\sigma)\), every restriction in
\[
 \mathcal P_\gamma
 =\{(q',x):q'\in R_n\cap\overline{\mathcal Q}\cap\overline\omega,
       \ A_\gamma x=b_\gamma+B_\gamma q',\ x\geq0\}
 \tag{2}
\]
is affine.  Indeed, \(R_n\) is the intersection of the two arm simplices and the finitely many closed interval restrictions in \(\texttt{def:observed-confidence-region}\); \(\overline{\mathcal Q}\) is the simplex intersection with the linear nonnegativity restrictions in \(\texttt{def:compatible-laws}\) and \(\texttt{def:observed-balances}\); and \(\overline\omega\) consists of weak affine comparisons between propensities and fixed grid values.  The remaining restrictions in (2) are precisely those of \(\texttt{def:joint-inversion-polytopes}\) and \(\texttt{def:allocation-fiber}\).  Thus \(\mathcal P_\gamma\) is a polyhedron.  It is bounded: the bin-balance rows in \(A_\gamma x=b_\gamma+B_\gamma q'\) imply
\[
 \sum_s x_{bs}=|I_b|\quad\hbox{for every bin }b,
 \qquad
 \sum_{b,s}x_{bs}=\sum_b|I_b|=1,
 \tag{3}
\]
while \(x\geq0\), and \(q'\) lies in a product of simplices.  Hence \(\mathcal P_\gamma\) is a compact polytope, including when some retained tied bin has width zero.  Its image under \((q',x)\mapsto C_\gamma x\) is therefore a compact polytope.

For completeness, every \(q'\in\overline{\mathcal Q}\) has at least one feasible path allocation.  To see this without adding a support assumption, put \(m=p_1-p_0\).  The observable balances give
\[
 \sum_i c_{0i}=m=\sum_j c_{1j},\qquad c_{0i}\geq0,\quad c_{1j}\geq0.
 \tag{4}
\]
Partition \([0,m]\) consecutively into intervals of lengths \((c_{01},\ldots,c_{0K})\) and, separately, into intervals of lengths \((c_{11},\ldots,c_{1K})\).  Let \(g_{ij}\) be the length of the intersection of the \(i\)-th interval of the first partition and the \(j\)-th interval of the second.  Then
\[
 g_{ij}\geq0,qquad \sum_jg_{ij}=c_{0i},qquad \sum_ig_{ij}=c_{1j}.
 \tag{5}
\]
The positive cells of \(g\) are coordinatewise ordered: if \(i<i'\) and both \(g_{ij}\) and \(g_{i'j'}\) are positive, then \(j\leq j'\).  Their support is therefore contained in a maximal monotone lattice path \(\sigma\).  Such a path visits every row and every column.  For each column \(j\), choose a visited node \((i(j),j)\) and put the always-treated mass \(a_{1j}\) there; for each row \(i\), choose a visited node \((i,j(i))\) and put the never-treated mass \(n_{0i}\) there; put the complier mass \(g_{ij}\) at its path node.  Split each of these three path-node mass vectors over the policy bins inside its principal-stratum interval in proportion to bin width.  For example, when \(m>0\), the complier-bin allocation is
\[
 x_{bs}=\frac{|I_b|}{m}h^{c}_s,
 \qquad b\subseteq(p_0,p_1],
 \tag{6}
\]
where \(h^c_s\) is the mass from (5) at node \(s\); the analogous formulas use denominators \(p_0\) and \(1-p_1\) in the always-treated and never-treated intervals.  If one of these denominators is zero, (4) and the simplex identities make the associated masses and bin widths zero, and the corresponding allocation is set to zero.  Equations (3)--(6) verify every row and outcome-balance equation in \(\texttt{def:allocation-fiber}\).  Thus some \(X_{(\omega,\sigma)}(q')\) is nonempty for every order closure containing \(q'\).  By the one-law attainability part of \(\texttt{thm:path-representation}\), the resulting image is generated by a maintained law.  In particular, all sharp sets parameterized below are nonempty compact sets; no complier full-support condition was used in (4)--(6).

Now apply \(\texttt{thm:path-representation}\) separately at each \(q'\in\mathcal Q_n\) and take the union.  The definition \(\texttt{def:outer-inversion}\) yields exactly
\[
 C_n
 =\bigcup_{\gamma\in\Gamma_K}
   \{C_\gamma x:(q',x)\in\mathcal P_\gamma\}.
 \tag{7}
\]
This is a finite-union extended formulation, and (2) shows that it contains no relaxation.  At a tie, repeated order labels merely duplicate an image and hence do not change (7).  The same theorem, with the law variable held fixed across all components, gives
\[
 \mathcal H_n
 =\left\{
   \bigcup_{\substack{\gamma=(\omega,\sigma)\in\Gamma_K\\q'\in\overline\omega}}
   \{C_\gamma x:A_\gamma x=b_\gamma+B_\gamma q',\ x\geq0\}
   :q'\in\mathcal Q_n
 \right\}.
 \tag{8}
\]
Thus (8) is an exact common-parameter finite-union representation of every member of \(\mathcal H_n\).  The common \(q'\) in (8) is indispensable: independently varying it across components would describe unions assembled from different observed laws, which are not members of \(\mathcal H_n\).

It remains to prove that the two metric radii are exactly computable.  Consider two nonempty compact sets represented as finite unions of the images of affine systems in (2).  Their Hausdorff distance is at most \(r\) if and only if both directed conditions hold:
\[
 \forall a\in A\ \exists b\in B:\ -r\leq a_\ell-b_\ell\leq r
       \quad(\ell=1,\ldots,L),
 \tag{9}
\]
\[
 \forall b\in B\ \exists a\in A:\ -r\leq a_\ell-b_\ell\leq r
       \quad(\ell=1,\ldots,L).
 \tag{10}
\]
This equivalence follows directly from the definition of Hausdorff distance induced by \(\|\cdot\|_\infty\): (9) and (10) say that each of the two directed distances is at most \(r\).  Replace membership in \(A\) and \(B\) by the finite disjunction over their component labels and introduce the associated variables \((q',x)\) from (2).  Replace \(a\) and \(b\) by \(C_\gamma x\).  Then (9)--(10) become a first-order formula whose atomic predicates are affine equalities or inequalities.  Universal membership assertions are written as implications, so empty individual components cause no difficulty, while the nonemptiness established above ensures that Hausdorff distance itself is well defined.

More explicitly, let \(\Phi(q^a,q^b,r)\) denote the resulting linear-real formula for
\(d_H(S(q^a),S(q^b))\leq r\), and let \(\Phi_{\mathrm{out}}(q,r)\) be the analogous formula for \(d_H(C_n,S(q))\leq r\), using (7) for membership in \(C_n\).  From \(\texttt{def:inversion-radii}\),
\[
 r_n^{\mathrm{out}}
 =\inf\{r\geq0:\ \forall q\in\mathcal Q_n,\ \Phi_{\mathrm{out}}(q,r)\},
 \tag{11}
\]
because the set on the right consists exactly of the upper bounds on
\(\{d_H(C_n,S(q)):q\in\mathcal Q_n\}\).  Likewise,
\[
 r_n^{\mathrm{set}}
 =\inf\{r\geq0:\ \exists q^a\in\mathcal Q_n\ \forall q^b\in\mathcal Q_n,
                      \ \Phi(q^a,q^b,r)\}.
 \tag{12}
\]
Equation (12) is exactly the infimum, over the center \(S(q^a)\in\mathcal H_n\), of the supremum of its distances to all \(S(q^b)\in\mathcal H_n\).  These equalities concern infima of definable upper-bound sets and therefore remain valid even if a displayed infimum or supremum is not attained.

We next give a terminating elimination argument and its operation bound.  By the definitions of \(v_{\mathrm{alg}}\) and \(m_{\mathrm{alg}}\) in \(\texttt{def:finite-inversion-algorithm}\), after expanding the two directions (9)--(10), either radius formula has at most
\[
 v\leq v_{\mathrm{alg}}
 \quad\text{quantified real variables and}\quad
 M_0\leq 2m_{\mathrm{alg}}|\Gamma_K|
 \quad\text{scalar affine atoms}.
 \tag{13}
\]
The Clopper--Pearson endpoints and all entries of the component arrays are inputs and hence are constants in these formulas.

To eliminate one existentially quantified variable \(w\), write each atom involving it as
\[
 a_iw+\ell_i(\eta)\ \mathrel{\bowtie_i}\ 0,
 \tag{14}
\]
where \(\eta\) contains the remaining variables, \(a_i\) is an input constant, and \(\bowtie_i\) is one of the order or equality relations.  A comparison in the ordered field decides whether \(a_i=0\).  For \(a_i\neq0\), the only value of \(w\) where the truth of this atom can change is
\[
 \rho_i(\eta)=-\ell_i(\eta)/a_i.
 \tag{15}
\]
For fixed \(\eta\), the Boolean formula is constant on every open interval between consecutive distinct values in (15), and boundary truth is tested at the values themselves.  Hence \(\exists w\) is eliminated by the finite disjunction obtained by testing: every \(\rho_i\); the midpoint of every pair that is consecutive under the corresponding affine ordering conditions; one point below a candidate minimum; and one point above a candidate maximum.  The assertions that a candidate pair is consecutive or that a candidate is extreme are themselves affine comparisons among the \(\rho_i\)'s.  This construction handles strict and weak atoms separately at the boundary and uses at most a fixed constant times \(M^3\) atoms and ordered-field operations when the incoming formula has \(M\) atoms.  If (14) has no nonzero coefficient, the formula is independent of \(w\).  A universal quantifier is eliminated by applying the same construction to the negated formula.

Iterating this explicit construction for the \(v\) variables in (13) gives, for a numerical constant \(c\), a recurrence of the form
\[
 M_{j+1}\leq cM_j^3,
 \qquad
 M_v\leq M_0^{,3^{v}O(1)}=M_0^{,2^{O(v)}}.
 \tag{16}
\]
The number of arithmetic and order operations obeys the same bound.  Combining (13), (16), and (1) proves the claimed unit-cost bound
\[
 \bigl(2m_{\mathrm{alg}}|\Gamma_K|\bigr)^{2^{O(v_{\mathrm{alg}})}}
 \quad\text{with}\quad
 |\Gamma_K|\leq(2L+1)^2\binom{2K-2}{K-1}.
 \tag{17}
\]
After elimination, each predicate in (11) or (12) is a Boolean combination of affine comparisons in the single variable \(r\).  Its truth set is therefore a finite union of intervals and points.  Comparing its finitely many affine boundary values and testing the intervening cells isolates the exact boundary giving the infimum in (11) or (12), whether the boundary is included or excluded.  This proves both termination and exactness of the two radii on the nonempty branch.

Finally suppose \(\mathcal Q_n=\varnothing\).  Emptiness is itself decided by the same finite affine elimination.  The algorithm then invokes exactly the symbolic fallback in \(\texttt{def:outer-inversion}\) and \(\texttt{def:inversion-radii}\).  Put \(\Delta=y_K-y_1>0\), let
\[
 X=[y_1,y_K]^L,qquad
 x^-=(y_1,\ldots,y_1),\qquad x^+=(y_K,\ldots,y_K).
\]
For every nonempty compact \(A\subseteq X\), \(d_H(X,A)\leq\Delta\), while
\(d_H(X,\{x^-\})=\|x^+-x^-\|_\infty=\Delta\).  Therefore the fallback outer radius is exactly
\[
 \sup_{\varnothing\neq A\subseteq X\ \mathrm{compact}}d_H(X,A)=\Delta.
 \tag{18}
\]
Let \(x^0=((y_1+y_K)/2,\ldots,(y_1+y_K)/2)\).  Every point of \(X\) is within \(\Delta/2\) of \(x^0\), so
\[
 \sup_{\varnothing\neq B\subseteq X\ \mathrm{compact}}d_H(\{x^0\},B)\leq\Delta/2.
 \tag{19}
\]
Conversely, for any nonempty compact \(A\subseteq X\), the triangle inequality gives
\[
 \Delta=d_H(\{x^-\},\{x^+\})
 \leq d_H(\{x^-\},A)+d_H(A,\{x^+\}),
 \tag{20}
\]
so at least one of the two distances on the right is at least \(\Delta/2\).  Taking the supremum over fallback members and then the infimum over \(A\) proves from (19)--(20) that the fallback set-of-sets radius is exactly \(\Delta/2\).  This branch returns the prescribed collection symbolically and makes no finite-polytope claim for its arbitrary compact members.  Equations (1)--(20) establish every assertion of the theorem.
